\documentclass[11pt]{article}
\usepackage{estudio}
\cuadernillo{2 --- Implementación}

\begin{document}

\portada{2}{Implementación}{%
  Diapositivas 5 a 7 \textbullet\ expositor B \textbullet\ 1:37\\[4pt]
  Arquitectura del motor y Cell Index Method}

\tableofcontents
\clearpage

% ===========================================================================
\section{Qué se cuenta acá y qué no}

La guía de la cátedra (punto 2.2) delimita esta sección con precisión:

\begin{quote}
``Detalle de cómo se traduce el modelo matemático al modelo computacional,
describiendo arquitectura del código implementado, pseudocódigo, diagrama UML,
etc. \textbf{Solo considerar el motor de simulación propiamente dicho, dejando
fuera de esta descripción cómo se implementa el post-proceso, o el formato de
los archivos input/output}, etc.''
\end{quote}

\ojo{No hables del formato de los archivos de texto, ni de \texttt{sweep.py},
ni de \texttt{analyze.py}, ni de cómo se hacen las animaciones. Está
explícitamente excluido. Si te preguntan por el output, contestás en el turno de
preguntas (está en el cuadernillo 3), pero no lo pongas en el guion.}

Tenés dos diapositivas de contenido y 1:37 en total. Es el bloque más corto:
sé quirúrgico.

% ===========================================================================
\section{Guion: diapositivas 5 a 7}

\diapo{5}{Divisoria: Implementación}{2 s}{B}

Solo el título de la sección. Es tu señal de entrada.

\diapo{6}{Arquitectura del motor}{50 s}{B}

A la izquierda el diagrama de flujo del paso temporal; a la derecha cuatro
puntos. Recorré el diagrama de arriba hacia abajo y cerrá con la sincronía.

\begin{quote}
\itshape
``El motor está escrito en C++20, sin dependencias externas. La clase
\texttt{Simulation} concentra todo el estado: las partículas, la grilla y el
generador pseudoaleatorio. La partícula guarda solamente su posición y su
ángulo: la velocidad no se almacena, se deriva del ángulo, porque el módulo es
constante. Cada paso hace tres cosas: primero reconstruye la grilla del Cell
Index Method para obtener la lista de vecinos; después calcula las
orientaciones nuevas con la regla de alineación, que es lo único que distingue
Vicsek del votante; y por último integra las posiciones aplicando las
condiciones periódicas. Los dos modelos comparten absolutamente todo el código
salvo esa función de alineación. Y algo importante: la actualización es
sincrónica, las orientaciones nuevas se escriben en un buffer aparte, así que
ninguna partícula ve el ángulo nuevo de otra dentro del mismo paso.''
\end{quote}

\clave{Los dos puntos que no podés dejar de decir: \textbf{(1)} los dos modelos
comparten todo el motor menos una función, y \textbf{(2)} la actualización es
sincrónica gracias a un buffer auxiliar. El primero es lo que hace legítima la
comparación de la sección de resultados; el segundo es lo que hace que esto sea
un autómata celular bien planteado.}

\diapo{7}{Búsqueda de vecinos: Cell Index Method}{45 s}{B}

A la izquierda el esquema de la grilla con la celda central resaltada y la
condición $L/M > r_c$; a la derecha cuatro puntos.

\begin{quote}
\itshape
``Para la búsqueda de vecinos reutilizamos el Cell Index Method del TP1.
Dividimos la caja en una grilla de $M$ por $M$ celdas, con el lado de celda
elegido de forma que sea mayor que el radio de interacción. Esa condición
garantiza que todos los vecinos posibles de una partícula estén en su propia
celda o en las ocho adyacentes, así que solo comparamos contra ese vecindario en
lugar de contra las $N$ partículas. Con las condiciones periódicas, la última
fila de celdas es adyacente a la primera, así que los índices se cierran sobre
el toro. El costo pasa de orden $N^2$ a orden $N$. Y como las partículas se
mueven, la grilla se reconstruye en cada paso.''
\end{quote}

\clave{La frase clave es \emph{``el lado de celda tiene que ser mayor que el
radio de interacción''}: es la condición que hace correcto al método. Si el
lado fuera menor, habría vecinos afuera del bloque de 3$\times$3 y se
perderían.}

\ojo{La diapositiva dice ``sólo se comparan las 9 celdas vecinas''. Eso es
correcto \emph{conceptualmente}, pero el código usa la optimización de medio
vecindario: recorre 5 celdas por celda y registra cada par encontrado en las
listas de las \emph{dos} partículas. El resultado es idéntico al de recorrer
las 9, con la mitad del trabajo. Si te preguntan, tenés que saberlo (ver P6).}

Cerrá pasándole la palabra a C:

\begin{quote}
\itshape
``Con el motor descripto, [C] pasa al sistema concreto que simulamos.''
\end{quote}

% ===========================================================================
\section{El fondo: cada decisión de implementación y su justificación}

Esta es la sección donde más pueden pinchar, porque cada línea de código es una
decisión que hay que poder defender. Van ordenadas por probabilidad de que las
pregunten.

\subsection{Por qué $M = 9$ y no $M = 10$}

El número de celdas por lado lo calcula \texttt{maxValidGridM(L, rc)}: devuelve
el mayor entero $M$ tal que el lado de celda $L/M$ sea \emph{estrictamente}
mayor que $r_c$.

\[
  M = \max(1,\ m), \qquad
  m = \begin{cases}
    \lfloor L/r_c \rfloor - 1, & \text{si } L/r_c \text{ es entero,} \\
    \lfloor L/r_c \rfloor,     & \text{en caso contrario.}
  \end{cases}
\]

Con $L = 10$ y $r_c = 1$ el cociente da exactamente 10, que es entero, así que
se resta uno y queda $\mathbf{M = 9}$, con lado de celda
$10/9 \approx 1{,}11 > 1$.

\dato{Si tomáramos $M = 10$, el lado de celda sería exactamente $r_c = 1$. Dos
partículas separadas por una distancia apenas menor que $r_c$ ---o sea,
vecinas--- podrían caer en celdas \emph{no} adyacentes, y el método las
perdería. La desigualdad tiene que ser estricta, y por eso se resta uno cuando
el cociente es entero. Buscar más celdas de las válidas no acelera nada: rompe
el método.}

\subsection{Por qué la grilla se reconstruye en cada paso}

Porque las partículas se mueven. En un autómata celular de red el vecindario es
fijo y se calcula una vez; acá, después de integrar las posiciones, una
partícula puede haber cambiado de celda y su conjunto de vecinos es otro. No
reconstruir significaría usar una lista de adyacencia desactualizada.

Ahora bien, \emph{reconstruir} no significa \emph{reasignar memoria}. La clase
\texttt{Grid} es un objeto persistente: los vectores de celdas y de listas de
vecinos se \emph{vacían} y se vuelven a llenar en cada llamada, sin liberar ni
volver a reservar. Eso importa porque el barrido paramétrico completo hace del
orden de $10^6$ reconstrucciones.

\teoria{Esa es la diferencia principal con el \texttt{computeCIM} del TP1, que
creaba y destruía sus buffers en cada llamada. Es también parte de por qué en el
benchmark (diapositiva 25) el CIM de TP2 sale \emph{más rápido} que el de TP1
para el mismo $N$.}

\subsection{Por qué el costo es $\mathcal{O}(N)$ y no $\mathcal{O}(N^2)$}

Con fuerza bruta se comparan todos los pares: $N(N-1)/2$ comparaciones. Con el
CIM, cada partícula solo se compara contra las que están en su celda y en las
adyacentes. Si la densidad se mantiene constante, el número esperado de
partículas por celda es constante, así que cada partícula hace un número
\emph{constante} de comparaciones y el total escala como $N$.

\ojo{La linealidad vale \textbf{a densidad constante}. En nuestro benchmark
aumentamos $N$ con la caja fija en $L=10$, así que la densidad crece con $N$ y
hay más partículas por celda: el escalado medido resulta superlineal. Esto lo
explica C en la diapositiva 25, pero conviene que lo sepas porque te lo pueden
preguntar acá.}

\subsection{Las dos condiciones periódicas, que son distintas}

La periodicidad interviene en dos lugares separados, y conviene no confundirlos:

\begin{enumerate}
  \item \textbf{En la indexación de celdas.} Los índices de celdas vecinas se
    cierran sobre el toro con $\mathrm{wrap}(i, M) = ((i \bmod M) + M) \bmod M$,
    de modo que la última fila de celdas es adyacente a la primera.
  \item \textbf{En el cálculo de la distancia.} El predicado de vecindad aplica
    la \emph{convención de imagen mínima}:
    $\Delta = d - L\,\mathrm{round}(d/L)$ a cada componente, antes de comparar
    contra $r_c^2$. Sin eso, dos partículas a ambos lados del borde darían una
    distancia enorme en lugar de la correcta.
\end{enumerate}

El constructor de \texttt{Grid} exige además $L/2 > r_c$. Esa condición
garantiza que la imagen mínima sea \emph{única}: si $r_c$ fuera mayor que media
caja, una partícula podría ser vecina de otra por dos caminos distintos
alrededor del toro y la convención dejaría de estar bien definida. Con $L=10$ y
$r_c=1$ se cumple holgadamente.

\subsection{Medio vecindario: 5 celdas en vez de 9}

Cuando la grilla tiene $M \geq 3$, el código recorre solo 5 de las 9 celdas del
vecindario: la propia, y las de arriba, derecha, arriba-derecha y
arriba-izquierda. Al encontrar un par $(i,j)$ lo registra en las listas de
\emph{ambas} partículas.

El resultado es idéntico al de recorrer las 9 ---cada par se descubre
exactamente una vez, desde uno de los dos lados--- pero con la mitad de las
comparaciones. Es la misma optimización que ya traía el CIM del TP1.

\teoria{La condición $M \geq 3$ no es caprichosa. Con $M < 3$ y contorno
periódico, el \textit{wrap} hace que una misma celda aparezca dos veces en el
stencil de medio vecindario y el par se registraría duplicado; por eso el
código cae al vecindario completo de 9 celdas en ese caso. Con $M = 9$ estamos
siempre en el camino optimizado.}

\subsection{La actualización sincrónica, en tres pasadas}

El método \texttt{step()} hace exactamente esto:

\begin{enumerate}
  \item \textbf{Reconstruye la grilla} sobre las posiciones del instante $t$.
  \item \textbf{Calcula todas las orientaciones nuevas} leyendo únicamente el
    estado viejo, y las deposita en un buffer aparte, \texttt{thetaNew\_}. El
    campo \texttt{theta} de las partículas no se escribe en ningún punto de este
    bucle.
  \item \textbf{Integra las posiciones} y recién ahí confirma los ángulos
    nuevos.
\end{enumerate}

\clave{Si no hubiera buffer auxiliar y se escribiera $\theta_i$ en el mismo
bucle, la partícula $i+1$ vería el ángulo \emph{nuevo} de la partícula $i$
dentro del mismo paso. Eso sería una actualización \emph{asincrónica}, el
resultado dependería del orden de recorrido de las partículas ---que es
arbitrario--- y dejaría de ser un autómata celular bien definido. Es la razón de
ser del buffer.}

\subsection{La convención de la posición: se avanza con la velocidad vieja}

Este punto es sutil y conviene tenerlo claro porque se corrigió después de la
clase de consultas. La integración es

\[
  \vect{x}_i(t+\Delta t) = \vect{x}_i(t) + \vect{v}_i(t)\,\Delta t,
\]

con $\vect{v}_i(t)$ derivada de $\theta_i(t)$, o sea del ángulo \textbf{viejo},
tal como lo escribe la diapositiva 42 de la Teórica 2. En el código: la tercera
pasada calcula el desplazamiento con \texttt{p.theta} (todavía el viejo) y
\emph{después} asigna \texttt{p.theta = thetaNew\_[i]}.

\dato{La consecuencia física es que el ruido sorteado en el paso $t$ recién
afecta el desplazamiento del paso siguiente: la partícula primero avanza en la
dirección que traía y después gira. La convención opuesta (girar y después
avanzar) da un sistema sistemáticamente menos polarizado; la diferencia en $v_a$
llegaba a $0{,}10$ en la región de transición. Usamos la de la teórica.}

\subsection{El vecindario es autoinclusivo}

$\mathcal{N}_i$ contiene a la propia partícula $i$ además de sus vecinos. En
Vicsek, la suma de senos y cosenos arranca con el $\theta_i$ propio antes de
sumar los vecinos ---es la convención del paper original---. En el votante, el
sorteo se hace sobre $\{i\} \cup \mathcal{N}_i$, así que una partícula puede
``votarse a sí misma''.

Esto resuelve un caso de borde real: \textbf{¿qué hace una partícula que no
tiene ningún vecino dentro de $r_c$?} Con la convención autoinclusiva, conserva
su propia orientación más el ruido, y la regla queda bien definida. Sin ella,
el promedio circular sería $0/0$ y el sorteo del votante sería sobre un conjunto
vacío.

\dato{No es un caso marginal. En el estudio de clusters a densidades
subcríticas ($N = 11$ a $32$, con $\langle k \rangle < 1$) las partículas
aisladas son frecuentes, y la convención afecta una fracción no despreciable de
las actualizaciones.}

\ojo{En el código, el sorteo del votante es
\texttt{uniform\_int\_distribution<size\_t> pick(0, row.size())}: el índice
igual a \texttt{row.size()} representa a la propia partícula. Escribirlo como
\texttt{(0, row.size()-1)} sobre un vecindario vacío daría
\emph{comportamiento indefinido} por desbordamiento del entero sin signo. Está
hecho a propósito así.}

\subsection{Por qué \texttt{atan2} de las sumas y no un promedio de ángulos}

Si promediáramos los ángulos directamente, dos partículas a $+179^\circ$ y
$-179^\circ$ ---que apuntan casi en la misma dirección--- darían un promedio de
$0^\circ$, o sea la dirección opuesta. Es la discontinuidad de la variable
angular en $\pm\pi$.

La solución es el \emph{promedio circular}: se suman los senos y los cosenos por
separado ---o sea, se suman los vectores unitarios--- y se recupera el ángulo
con \texttt{atan2} de esas dos sumas. Eso es geométricamente correcto en toda la
circunferencia.

\subsection{Por qué se renormaliza el ángulo a $[-\pi, \pi]$}

Después de sumar el ruido, el código hace
$\theta \leftarrow \mathrm{atan2}(\sin\theta, \cos\theta)$, que devuelve el
mismo ángulo pero llevado al intervalo $[-\pi, \pi]$.

En Vicsek sería redundante, porque \texttt{atan2} ya devuelve un valor en rango.
Pero en el votante es \textbf{necesario}: la partícula copia \emph{textualmente}
el ángulo ya comprometido de otra, en lugar de recalcularlo. Sin la
renormalización, el ángulo acumularía una caminata aleatoria no acotada
---sumando ruido sobre ruido a lo largo de miles de pasos--- y crecería sin
límite en valor numérico, con la consiguiente pérdida de precisión en los senos
y cosenos.

\subsection{El generador pseudoaleatorio}

Es un \texttt{std::mt19937\_64} (Mersenne Twister de 64 bits) sembrado
explícitamente, nunca desde el reloj. El ruido de los dos modelos pasa por la
\emph{misma} función \texttt{addAngularNoise}, así que ambos usan exactamente la
misma distribución y el mismo flujo de números aleatorios. Cómo se derivan las
semillas del barrido lo cuenta C (cuadernillo 3).

\subsection{Compilación}

C++20, compilado con \texttt{-O2 -Wall -Wextra -pedantic}, sin dependencias
externas ni CMake: un \texttt{Makefile} plano. Hay además un binario de
autotest (\texttt{tp2\_test}) que verifica invariantes del CIM ---simetría de la
lista de adyacencia, ausencia de auto-vecinos y de duplicados--- contra una
fuerza bruta de referencia.

\dato{El \texttt{-O2} no es cosmético: las mediciones de tiempo del inciso (g)
carecen de sentido sin optimización.}

% ===========================================================================
\section{Preguntas y respuestas}

\subsection{Sobre el Cell Index Method}

\pregunta{¿Por qué $M = 9$ y no $M = 10$, si $L = 10$ y $r_c = 1$?}
\respuesta{Porque el lado de celda tiene que ser \emph{estrictamente} mayor que
$r_c$. Con $M = 10$ el lado sería exactamente $1$, y dos partículas separadas
por una distancia apenas menor que $r_c$ podrían caer en celdas no adyacentes y
perderse como vecinas. Por eso la fórmula resta uno cuando $L/r_c$ da entero:
queda $M = 9$ y el lado de celda es $10/9 \approx 1{,}11$.}

\pregunta{¿Qué pasaría si usaran un $M$ más grande? ¿No sería más rápido?}
\respuesta{No: sería incorrecto. El método deja de encontrar todos los vecinos.
El constructor de la grilla directamente lanza una excepción si le pasás un $M$
mayor que el máximo válido.}

\pregunta{¿Por qué reconstruyen la grilla en cada paso? ¿No es caro?}
\respuesta{Es necesario, porque las partículas se mueven y cambian de celda. Lo
que sí evitamos es el costo de reservar memoria: la grilla es un objeto
persistente cuyos vectores se vacían y se vuelven a llenar, sin liberar ni
volver a reservar. Sobre el orden de $10^6$ reconstrucciones que hace el barrido
completo, eso es lo que importa.}

\pregunta{Dicen ``9 celdas'', pero ¿el código recorre las 9?}
\respuesta{No, recorre 5 y registra cada par en las listas de las dos
partículas. El resultado es idéntico al de recorrer las 9 ---cada par se
descubre una sola vez, desde uno de los dos lados--- con la mitad de las
comparaciones. Es la optimización de medio vecindario, la misma del TP1.
Conceptualmente el vecindario sigue siendo el bloque de $3 \times 3$.}

\pregunta{¿Cómo manejan las condiciones periódicas en el CIM?}
\respuesta{En dos lugares distintos. En la indexación, los índices de celda se
cierran sobre el toro con una operación de módulo, así que la última fila es
adyacente a la primera. Y en la distancia, aplicamos la convención de imagen
mínima antes de comparar contra el radio. Además exigimos $L/2 > r_c$ para que
la imagen mínima sea única; con $L = 10$ y $r_c = 1$ se cumple de sobra.}

\pregunta{¿El CIM es realmente $\mathcal{O}(N)$?}
\respuesta{A densidad constante, sí: el número esperado de partículas por celda
es constante, así que cada partícula hace un número constante de comparaciones.
En nuestro benchmark aumentamos $N$ con la caja fija, así que la densidad crece
y el escalado medido da superlineal. Eso está discutido en la diapositiva 25.}

\pregunta{¿Verificaron que el CIM da lo mismo que la fuerza bruta?}
\respuesta{Sí, hay un binario de autotest que compara la lista de adyacencia
del CIM contra una fuerza bruta de referencia y verifica los invariantes:
simetría, que ninguna partícula sea vecina de sí misma y que no haya
duplicados.}

\subsection{Sobre la actualización}

\pregunta{¿Por qué la actualización es sincrónica? ¿Qué cambiaría si no lo
fuera?}
\respuesta{Las orientaciones nuevas se escriben en un buffer aparte y recién se
confirman al final del paso, así que ninguna partícula ve el ángulo actualizado
de otra dentro del mismo paso. Si fuera asincrónica, el resultado dependería del
orden en que se recorren las partículas, que es arbitrario, y el modelo dejaría
de estar bien definido como autómata celular.}

\pregunta{¿La posición se actualiza con la velocidad vieja o con la nueva?}
\respuesta{Con la vieja: $\vect{x}(t+\Delta t) = \vect{x}(t) +
\vect{v}(t)\,\Delta t$, tal como está en la diapositiva 42 de la Teórica 2. La
partícula primero avanza en la dirección que traía y después gira, así que el
ruido sorteado en un paso recién afecta el desplazamiento del paso siguiente.}

\pregunta{¿Qué hace una partícula que se queda sin vecinos?}
\respuesta{El vecindario es autoinclusivo: contiene a la propia partícula. En
Vicsek, el promedio circular arranca sumando su propio seno y coseno ---es la
convención del paper original---, así que una partícula aislada conserva su
orientación más el ruido. En el votante, el sorteo se hace sobre el conjunto que
incluye a la propia partícula, así que se ``vota a sí misma''. En los dos casos
la regla queda bien definida. No es un caso marginal: en el estudio de clusters
a densidades subcríticas las partículas aisladas son frecuentes.}

\pregunta{¿Por qué usan \texttt{atan2} de las sumas de senos y cosenos en vez de
promediar los ángulos?}
\respuesta{Porque promediar ángulos directamente rompe en $\pm\pi$: dos
partículas a $+179^\circ$ y $-179^\circ$ apuntan casi igual, pero su promedio
aritmético da $0^\circ$, la dirección opuesta. Sumando los senos y cosenos por
separado estamos sumando vectores unitarios, que es lo geométricamente
correcto, y recuperamos el ángulo con \texttt{atan2}.}

\pregunta{¿Por qué renormalizan el ángulo a $[-\pi,\pi]$?}
\respuesta{Por el modelo de votante. Ahí la partícula copia textualmente un
ángulo ya comprometido en vez de recalcularlo con \texttt{atan2}, así que sin
renormalizar el ángulo acumularía una caminata aleatoria no acotada a lo largo
de miles de pasos. En Vicsek sería redundante porque \texttt{atan2} ya devuelve
un valor en rango, pero usamos la misma función en los dos para que el
tratamiento del ruido sea idéntico.}

\subsection{Sobre el código y la comparación entre modelos}

\pregunta{¿Cuánto código comparten los dos modelos?}
\respuesta{Todo salvo una función. Hay un \texttt{enum} que selecciona entre
\texttt{circularMeanHeading} y \texttt{voterHeading}; el resto ---la grilla, el
ruido, la integración, las condiciones periódicas, los observables--- es el
mismo código. El ruido de los dos pasa por la misma función, con la misma
distribución y el mismo flujo de números aleatorios. Eso es lo que hace que la
comparación mida la diferencia entre las reglas y no diferencias de
implementación.}

\pregunta{¿Por qué la partícula no guarda la velocidad?}
\respuesta{Porque el módulo es constante por construcción: la velocidad queda
determinada por el ángulo. Guardar las dos componentes sería estado redundante
que habría que mantener consistente. Se derivan cuando hacen falta, con seno y
coseno.}

\pregunta{¿Qué generador de números aleatorios usan?}
\respuesta{Mersenne Twister de 64 bits, sembrado explícitamente, nunca desde el
reloj. Todo el barrido es reproducible bit a bit.}

\pregunta{¿Por qué C++ y no Python para el motor?}
\respuesta{Por el volumen del barrido: son 468 puntos de simulación por 10
semillas, con corridas de 2000 pasos ---y de 10\,000 en las densidades
subcríticas---, del orden de $10^6$ reconstrucciones de grilla. En Python no
cerraba. El análisis y los gráficos sí están en Python, sobre los archivos de
texto que escribe el motor.}

\pregunta{¿Reutilizaron el código del TP1?}
\respuesta{Sí, el Cell Index Method. Lo reorganizamos en una clase
\texttt{Grid} persistente en vez de la función que reasignaba sus buffers en
cada llamada, y le sacamos el radio de las partículas, porque acá son
puntuales. La lógica de indexación, el \textit{wrap} periódico y la imagen
mínima son las mismas.}

\end{document}
